
%%%%%%%%%%%%%%%%%% 插图环境 %%%%%%%%%%%%%%%%%%%

\chapter{插图环境}

\section{图的使用}

XeLaTeX 环境下可以插入 PDF、PNG、JPEG、BMP 等格式的图片 (EPS 图片可先用 \texttt{epstopdf} 转换为 PDF), 也可以用绘图包直接在 \LaTeX 中绘制图形, 如 tikz (示例见第二章的算法流程图). 值得注意的是 figure 环境是一个浮动体\index{浮动体}环境, LaTeX 不总是把浮动体放在你想要的地方, 但是 LaTeX 总是保证浮动体的相对顺序, 所以对图片 \verb|\label| 和 \verb|\ref| 的交叉引用就显得尤为重要。

\section{插图示例}

插入一个图形并居中放置, 如图~\ref{fig:sinx}.
\begin{figure}[htp!]
  \centering
  \includegraphics[width=7.2cm]{image1}
  \caption{函数 $y=\sin(x)$ 的图像}\label{fig:sinx}
\end{figure}

两个图左右并排放置, 共用一个标题, 如图~\ref{fig:image}.
\begin{figure}[htp!]
  \centering
  \includegraphics[width=0.45\linewidth]{image1}
  \hfill
  \includegraphics[width=0.45\linewidth]{image2}
  \caption{左: 图一的描述;~ 右:图二的描述}
  \label{fig:image}
\end{figure}

\clearpage
使用 minipage 排版并排插图, 每个图都有单独的标题. 通过 \verb|autoref| 引用图片: \autoref{fig:image1} 与 \autoref{fig:image2}.
\begin{figure}[htp!]
\begin{minipage}[t]{0.48\linewidth}
  \centering
  \includegraphics[width=0.9\linewidth]{image1}
  \caption{图一的描述}
  \label{fig:image1}
\end{minipage}
\hfill
\begin{minipage}[t]{0.48\linewidth}
\centering
   \includegraphics[width=0.9\linewidth]{image2}
   \caption{图二的描述}
   \label{fig:image2}
\end{minipage}
\end{figure}

使用 subcaption 宏包实现多图并排, 如图~\ref{fig:images}.
\begin{figure}[htp!]
  \centering
  \begin{subfigure}[b]{0.3\linewidth}
    \includegraphics[width=\linewidth]{image1}
    \caption{Subcaption A}
  \end{subfigure}\hfill
  \begin{subfigure}[b]{0.3\linewidth}
    \includegraphics[width=\linewidth]{image1}
    \caption{Subcaption B}
  \end{subfigure}\hfill
  \begin{subfigure}[b]{0.3\linewidth}
    \includegraphics[width=\linewidth]{image1}
    \caption{Subcaption C}
  \end{subfigure}

  \vspace{1em}
  \begin{subfigure}[b]{0.3\linewidth}
    \includegraphics[width=\linewidth]{image2}
    \caption{Subcaption D}
  \end{subfigure}\hfill
  \begin{subfigure}[b]{0.3\linewidth}
    \includegraphics[width=\linewidth]{image2}
    \caption{Subcaption E}
  \end{subfigure}\hfill
  \begin{subfigure}[b]{0.3\linewidth}
    \includegraphics[width=\linewidth]{image2}
    \caption{Subcaption F}
  \end{subfigure}
  \caption{六个图并排}
  \label{fig:images}
\end{figure}

